\clearpage\subsection{Анализ результатов}
Для одной заявки совмещение этапов невозможно. При стороне 2048 медианы составили 10,78 мс в последовательном режиме и 11,19 мс в конвейерном. При стороне 512 конвейер уступил во всех измеренных конфигурациях. В частности, для 16 заявок получено 30,23 и 52,91 мс соответственно. Для пакета из 32 малых изображений запись занимала больше времени, чем чтение и сегментация вместе. Потенциал ускорения за счёт перекрытия этих коротких операций в такой конфигурации ограничен, а накладные расходы на обмен заявками, планирование потоков и совместное использование ресурсов входят в замер.

При стороне 2048 и \(M=4,16,32\) медиана конвейерного режима оказалась ниже. Наибольшее отношение медиан \(1{,}41\) получено для 16 заявок: время уменьшилось со 159,12 до 112,63 мс. Для 32 заявок получено ускорение \(1{,}15\), а межквартильные интервалы двух режимов существенно перекрылись. Поэтому выигрыш для этого размера пакета описывается по медианам и зависит от условий запуска.

При пакетах из 32 больших изображений самым длительным этапом конвейера по измеренным длительностям была сегментация (7,60 мс против 3,18 мс для чтения и 4,58 мс для записи). Медианная задержка заявки составила 65,75 мс, поскольку включает ожидание в очередях. Таким образом, ускорение обработки всего пакета может сопровождаться увеличением задержки отдельной заявки.

\structuralpart{ЗАКЛЮЧЕНИЕ}
Разработана трёхступенчатая обработка заявок с чтением файла, сегментацией методом Оцу и сохранением отдельной маски. Конвейер передаёт заявки через ограниченные очереди и сохраняет порядок обработки на каждом этапе. Результаты последовательного и конвейерного режимов совпали с независимым эталоном на всех корректных тестовых входах.

Временная диаграмма показала перекрытие операций над разными заявками. Для изображений \(2048\times2048\) и пакета из 16 заявок получено наибольшее медианное ускорение среди исследованных конфигураций, равное 1,41. Для изображений \(512\times512\) и для одиночных заявок в исследованных условиях последовательный режим имел меньшую медиану времени.
